% Options for packages loaded elsewhere
\PassOptionsToPackage{unicode}{hyperref}
\PassOptionsToPackage{hyphens}{url}
\documentclass[
]{article}
\usepackage{xcolor}
\usepackage{amsmath,amssymb}
\setcounter{secnumdepth}{-\maxdimen} % remove section numbering
\usepackage{iftex}
\ifPDFTeX
  \usepackage[T1]{fontenc}
  \usepackage[utf8]{inputenc}
  \usepackage{textcomp} % provide euro and other symbols
\else % if luatex or xetex
  \usepackage{unicode-math} % this also loads fontspec
  \defaultfontfeatures{Scale=MatchLowercase}
  \defaultfontfeatures[\rmfamily]{Ligatures=TeX,Scale=1}
\fi
\usepackage{lmodern}
\ifPDFTeX\else
  % xetex/luatex font selection
\fi
% Use upquote if available, for straight quotes in verbatim environments
\IfFileExists{upquote.sty}{\usepackage{upquote}}{}
\IfFileExists{microtype.sty}{% use microtype if available
  \usepackage[]{microtype}
  \UseMicrotypeSet[protrusion]{basicmath} % disable protrusion for tt fonts
}{}
\makeatletter
\@ifundefined{KOMAClassName}{% if non-KOMA class
  \IfFileExists{parskip.sty}{%
    \usepackage{parskip}
  }{% else
    \setlength{\parindent}{0pt}
    \setlength{\parskip}{6pt plus 2pt minus 1pt}}
}{% if KOMA class
  \KOMAoptions{parskip=half}}
\makeatother
\usepackage{longtable,booktabs,array}
\newcounter{none} % for unnumbered tables
\usepackage{calc} % for calculating minipage widths
% Correct order of tables after \paragraph or \subparagraph
\usepackage{etoolbox}
\makeatletter
\patchcmd\longtable{\par}{\if@noskipsec\mbox{}\fi\par}{}{}
\makeatother
% Allow footnotes in longtable head/foot
\IfFileExists{footnotehyper.sty}{\usepackage{footnotehyper}}{\usepackage{footnote}}
\makesavenoteenv{longtable}
\setlength{\emergencystretch}{3em} % prevent overfull lines
\providecommand{\tightlist}{%
  \setlength{\itemsep}{0pt}\setlength{\parskip}{0pt}}
\usepackage{hyperxmp}
\usepackage{bookmark}
\IfFileExists{xurl.sty}{\usepackage{xurl}}{} % add URL line breaks if available
\urlstyle{same}
\hypersetup{
  pdftitle={The server is not a structural requirement: identifying accidental architectural roles under journaled programs},
  pdfauthor={Alvaro Rivera},
  pdfkeywords={\xmpquote{server-role}, \xmpquote{journaled
programs}, \xmpquote{distributed systems}, \xmpquote{analytic
typology}, \xmpquote{event sourcing}, \xmpquote{actor
model}, \xmpquote{bootstrap}, \xmpquote{local-first}},
  hidelinks,
  pdfcreator={LaTeX via pandoc}}

\title{The server is not a structural requirement: identifying
accidental architectural roles under journaled programs}
\author{Alvaro Rivera}
\date{2026-05-16}

\begin{document}
\maketitle

\subsection{TL;DR}\label{tldr}

The role we call \emph{server} --- the entity that owns coordination,
accepts authoritative writes, and mediates between clients --- is
treated across the architecture literature as a structural component of
any non-trivial distributed system. Prior work has questioned the
server's privileges (local-first, federation), eliminated it as a design
goal (Hypercore, Holochain), or replaced it with a different topology
(peer-to-peer, blockchain). None of these positions removes the
server-role as a \emph{prerequisite for software construction itself}.

This paper applies Paper 6's \emph{infrastructural symptom} lens at the
role layer --- what Paper 6 §9 anticipates --- and names the role-level
expression as \emph{accidental category}. Under a class of systems in
which nodes replicate programs rather than data or state, the
server-role does not survive as a necessity. The bootstrap of a new
node, the last role that appeared to require a running service, can be
performed by an analog artifact: a printed QR code transferred by any
means and authorized by a human decision.

The lens is diagnostic, not prescriptive. It enables reading distributed
systems for whether the server-role exists in them as a structural
requirement or as a contingent artifact of the persistence model. An
instantiation --- a port of the Ordering bounded context of
\texttt{dotnet/eShop} to a journaled-program substrate, replicated
across three nodes that bootstrap through paper QR codes --- makes the
role-level reading operationally visible.

The contributions are two: the role-level extension of Paper 6's
typology (conceptual), and the working instantiation that makes it
inspectable (operational), together with the §6 observation that under
the same conditions cloud ecosystems pass from architectural
dependencies to library invocations.

\subsection{Formal statement}\label{formal-statement}

\textbf{Definition.} A \emph{server-role} is the architectural function
of accepting authoritative writes, owning coordination, mediating
between clients, and providing the address against which other parties
direct their traffic in a distributed system. The role is independent of
any specific machine, deployment topology, or cloud provider; it is the
function, not its incarnation.

A server-role reads as an \emph{accidental category} --- the role-level
expression of Paper 6's \emph{infrastructural symptom} --- in a class of
distributed systems S when both of two diagnostic indicators are
positive:

\begin{itemize}
\tightlist
\item
  \textbf{Contingency.} The role's necessity in S is traceable to a
  specific representational choice --- what the nodes of S replicate in
  order to share state. (This is Paper 6's \emph{compensation} indicator
  at the role layer.)
\item
  \textbf{Dissolution.} The role loses justification when that
  representational choice is replaced; nothing else about the problem
  domain, workload, or operational context needs to change for the role
  to become unnecessary. (This is Paper 6's \emph{dissolution} indicator
  unchanged.)
\end{itemize}

\textbf{Claim.} Under a class of systems in which nodes replicate
programs rather than data, state, or operations, both indicators read
positively for the server-role. The indicators are stated in §3; the
substrate that makes them read positively is exhibited in §4.

\textbf{Evidence.} Three operating nodes derived from a port of the
Ordering bounded context of \texttt{dotnet/eShop} to a journaled-program
substrate, bootstrapped by paper QR codes and authorized by human
decisions, execute the same domain behavior as the original microservice
across its operational lifecycle, without a node fulfilling the
server-role at any point in time. The original deployment topology
required at least one such role; the ported topology requires none.

\subsection{Claims this paper makes}\label{claims-this-paper-makes}

\begin{enumerate}
\def\labelenumi{\arabic{enumi}.}
\tightlist
\item
  Paper 6's \emph{infrastructural symptom} construct extends from
  infrastructural layers to architectural roles; the server-role reads
  as an accidental category --- the role-level expression of that
  construct --- under the same diagnostic typology (§3).
\item
  At the role layer, the same two diagnostic indicators of Paper 6 ---
  re-expressed as contingency and dissolution --- read the server-role
  as model-dependent or structural per use-case (§3).
\item
  Under a class of systems in which nodes replicate programs rather than
  data, both indicators read positively for the server-role: it
  disappears as a structural consequence, not as a design objective
  (§5).
\item
  The bootstrap of a new node --- the operation that prior literature
  treats as the residual case requiring a persistent server --- admits a
  presentation in which the information transferred crosses by analog
  means (F2) and the authorization is exercised by humans (F4), and the
  residual software discharge (F1, F3, F5) is bounded to the duration of
  the handshake, after which the issuer process exits (§4).
\item
  A working instantiation, constructed by porting the Ordering bounded
  context of \texttt{dotnet/eShop} to a journaled-program substrate and
  deploying it across three nodes bootstrapped by paper QR codes, makes
  the role-level reading operationally inspectable: the server-role
  dissolves across the operational lifecycle (§4, §5).
\item
  An unexpected consequence --- not anticipated by Paper 6 --- is that
  existing cloud and microservice ecosystems cease to be architectural
  dependencies and become reusable domain libraries within the program
  itself (§6).
\item
  The lens is diagnostic: it enables reading which architectural
  categories in any given system are structural and which are
  contingent, independent of the choice to act on that information (§7,
  §8).
\end{enumerate}

\subsection{1. Introduction}\label{introduction}

This paper extends the \emph{infrastructural symptom} construct of Paper
6 from infrastructural layers to architectural roles, an extension that
Paper 6 §9 anticipates. The genre is closer to \emph{theory for
analyzing} in the sense of Gregor (2006) than to design theory: the
construct itself is Paper 6's, applied at a different level of
abstraction. The novel material of this paper is the role-level
demonstration --- a working instantiation in which the server-role
dissolves --- and the operational observation that, under the same
conditions, cloud ecosystems pass from architectural dependencies to
libraries invoked by the program. The contribution is conceptual at the
extension level and operational at the demonstration level.

Production software is built around an assumption that is rarely
articulated as such: that a \emph{server-role} exists somewhere in the
system. The role is not the physical machine, nor the deployment
topology, nor any particular cloud provider. It is the architectural
function --- accepting authoritative writes, owning coordination,
mediating between clients, providing the address against which other
parties direct their traffic. Across textbooks, reference architectures,
and tooling defaults, the role is treated as a primitive of distributed
software: present by default, named when necessary, removed only at
great cost.

Prior literature has interrogated the server's \emph{privileges} --- its
asymmetric power over data, identity, or availability --- and proposed
designs that reduce those privileges. Local-first software argues that
data should belong to its user and survive without the cloud service.
Hypercore, Dat, and related peer-to-peer systems replicate append-only
logs without a privileged coordinator. Holochain inverts the
data-centric model so that the agent's source chain and a shared
distributed hash table together replace the central server. Federated
systems pluralize the server-role across cooperating hosts. Each of
these approaches treats the server-role as a \emph{target of design} ---
something to redesign, redistribute, or replace. None of them removes
the server-role as a \emph{prerequisite for software construction}. In
each case, the role's elimination is the explicit objective of the
system; the system is built to satisfy that objective.

This paper takes a different position. Under a class of systems in which
the substrate shared between nodes is the program (in the sense
established in Paper 2 §1.2: the pair of domain library and journal of
invocations) --- not the data, not the state, not even the operations
--- the server-role does not survive as a necessity. It is a category
whose referent becomes contingent. Server ceases to be a prerequisite
for making software. Bootstrap --- the last role that seemed to need a
persistent service --- can be discharged through a printed QR code as
its transfer medium, with the issuer's software footprint bounded to the
duration of the handshake.

The paper develops two claims. The first identifies the server-role as
an accidental category --- a role that exists contingently in
distributed software, as a consequence of choices about what nodes
replicate, and whose disappearance under different choices is a
structural consequence of those choices rather than a design goal. In
the class of systems where programs are journaled and replayed against
shared domain libraries, the role's elimination is what the substrate
produces, not what the engineer pursues. The second observes that the
displacement extends beyond the role itself: under the same conditions,
existing cloud and microservice ecosystems cease to function as
architectural dependencies and become reusable domain libraries that
programs can invoke at will. The cloud is not removed; its position in
the system is restated.

The instantiation used to demonstrate these claims --- a port of the
Ordering bounded context from Microsoft's \texttt{dotnet/eShop}
reference application, executed across three nodes that share a
journaled domain library and bootstrap each other through paper QR codes
--- is one of the paper's central contributions, not merely an existence
proof: what Paper 6's lens reads at the role level, this instantiation
makes operationally visible. The same lens admits other instantiations:
any system that satisfies the indicators of §3 would dissolve the
server-role in the same way. The choice of \texttt{dotnet/eShop} is
rhetorical --- Microsoft's widely-used reference application for
cloud-microservices patterns in the .NET ecosystem, whose original
deployment topology makes the contrast visible --- rather than essential
to the argument. The framework underlying the instantiation was not
designed to support this paper's claim. The property is identified
retrospectively, after a chain of unrelated architectural decisions
converged on a system in which the server-role had become contingent.
The work of this paper is, in that sense, a forensic reading of a class
of artifacts in which the property happens to hold --- consistent with
the analytic genre Gregor (2006) describes as \emph{theory for
analyzing}.

Section 2 situates the paper among prior work on local-first,
peer-to-peer, agent-centric, and federated systems, identifying the
dimension along which each work treats the server-role and the dimension
along which the present lens differs. Section 3 names \emph{accidental
category} as the role-level expression of Paper 6's
\emph{infrastructural symptom}, and re-expresses the two diagnostic
indicators --- contingency and dissolution --- at the role layer.
Section 4 presents the instantiation: the domain-library port, the
three-node deployment, the analog bootstrap. Section 5 examines the
first consequence: the dissolution of the server-role across the
operational lifecycle. Section 6 examines the second consequence --- the
one that prior literature does not anticipate: the conversion of
existing cloud ecosystems from architectural dependencies into reusable
domain libraries. Section 7 examines limitations, counter-arguments, and
the boundary of the lens's applicability. Section 8 concludes.

\subsection{2. Related work}\label{related-work}

Local-first software (Kleppmann et al., 2019) articulates seven ideals
for software that operates primarily on devices owned by their users
rather than on remote services. Replication is mediated by Conflict-free
Replicated Data Types (CRDTs), which allow concurrent edits to merge
without coordination. The position with respect to cloud infrastructure
is precise: cloud services are admissible, but the software must
continue to function in their absence. The title of the paper ---
\emph{Local-First Software: You Own Your Data, in spite of the Cloud}
--- captures the stance: the cloud is neither rejected nor required. The
server-role is reduced to one of several optional services, and its
elimination is the explicit objective of the architectural pattern.

Peer-to-peer append-only logs constitute a second line. Secure
Scuttlebutt (Tarr et al., 2019) replicates per-identity signed logs
across a gossip network without a central coordinator; the log, not the
data it carries, is the unit of synchronization. Hypercore (Holepunch,
n.d.) and the Dat protocol (Ogden et al., 2017) generalize the same
pattern, treating the append-only feed as a primitive replicated by a
distributed hash table. In each case the unit of replication is a log of
data --- events, document edits, file fragments --- and the application
logic that interprets the log lives in the consuming application. The
server's elimination is the design target of the network layer.

Holochain (Harris-Braun et al., 2018) makes the agent the source of
truth. Each agent maintains a local source chain --- an immutable,
signed sequence of its own transactions --- and contributes to a
validating distributed hash table. The architecture inverts the
data-centric assumption of blockchain rather than addressing the cloud
directly; the elimination of the server-role is a consequence of the
agent-centric inversion. The replicated artifacts are the agent's chain
and the validated DHT; application logic, written as zomes, is
interpreted locally by each agent.

Federated systems form a fourth line. Matrix federates an event graph
across cooperating homeservers; ActivityPub federates streams of
activities; the AT Protocol of Bluesky splits hosting across a Personal
Data Server, a Relay, and an AppView, with account portability as the
displacement mechanism. The server-role is preserved but pluralized:
every participating host fulfills the role for some subset of users. The
cloud is partitioned across cooperating providers rather than
centralized or rejected.

Adjacent work admits brief mention. Solid (Sambra et al., 2016) proposes
Personal Online Datastores that decouple identity from application
providers; the server-role moves from application owner to data host
without being eliminated. Earthstar (Earthstar Project, n.d.) pursues
local-first synchronization without a distributed hash table, sitting
between Kleppmann's CRDT model and Hypercore's gossip topology. Neither
treats the server-role as a category whose referent is contingent.

Each of these lines treats the server-role as a \emph{design target} ---
an architectural feature to redesign, redistribute, replace, or
pluralize. None of them names the property that the present paper
identifies: that the role is contingent on a choice about what nodes
replicate, and that under a different choice the role has no referent at
all. Table 1 summarizes the four families and the position of the
present construct.

\textbf{Table 1.} Five families of prior work and how each treats the
server-role. The bottom row identifies the construct of the present
paper; the columns describe what each family replicates between nodes,
where its application logic resides, why a privileged server is absent
(or pluralized), and how each family situates itself with respect to
cloud infrastructure.

{\def\LTcaptype{none} % do not increment counter
\begin{longtable}[]{@{}
  >{\raggedright\arraybackslash}p{(\linewidth - 8\tabcolsep) * \real{0.2000}}
  >{\raggedright\arraybackslash}p{(\linewidth - 8\tabcolsep) * \real{0.2000}}
  >{\raggedright\arraybackslash}p{(\linewidth - 8\tabcolsep) * \real{0.2000}}
  >{\raggedright\arraybackslash}p{(\linewidth - 8\tabcolsep) * \real{0.2000}}
  >{\raggedright\arraybackslash}p{(\linewidth - 8\tabcolsep) * \real{0.2000}}@{}}
\toprule\noalign{}
\begin{minipage}[b]{\linewidth}\raggedright
Family
\end{minipage} & \begin{minipage}[b]{\linewidth}\raggedright
What is replicated
\end{minipage} & \begin{minipage}[b]{\linewidth}\raggedright
Where the logic lives
\end{minipage} & \begin{minipage}[b]{\linewidth}\raggedright
How the server-role is treated
\end{minipage} & \begin{minipage}[b]{\linewidth}\raggedright
Relation to the cloud
\end{minipage} \\
\midrule\noalign{}
\endhead
\bottomrule\noalign{}
\endlastfoot
Local-first (Kleppmann) & CRDTs / data & In the application & Eliminated
as a design goal & Survival without it \\
Hypercore / Dat / SSB & Append-only logs & In the application &
Eliminated as a network goal & Replaced by peers \\
Holochain & Source chain + DHT & In the agent & Eliminated by
agent-centric inversion & Treated as an alternative architecture \\
Federated (Matrix / ActivityPub / AT Protocol) & Events / activities &
On federated servers & Pluralized across hosts & Fragmented across
providers \\
\textbf{The present construct} & \textbf{Programs} & \textbf{In the
journal of each node} & \textbf{Not a prerequisite (accidental
category)} & \textbf{Subsumed as a library} \\
\end{longtable}
}

The discriminator is the fourth column. The four families treat the
server-role as something to be designed away --- a residual feature
whose elimination is the goal. The present construct identifies the role
as contingent: under a representational choice in which nodes replicate
programs rather than data, state, or operations, the role's necessity
does not survive. The fifth column is downstream of the fourth: under
the present construct the cloud is neither evaded, replaced, nor
pluralized, but reused as a library within the program.

Beyond these four families that treat the server-role as a design
target, a fifth body of work shares direct lineage with the substrate of
§4 without making the server-role its design objective. Erlang/OTP
(Armstrong, 2003) is the longest-running production actor system; its
supervision trees and gen\_server abstraction shape the ecosystem from
which the substrate descends, though OTP nodes carry explicit names and
gen\_server processes carry explicit server-role at the §3 level.
Microsoft Orleans (Bernstein et al., 2014) introduces virtual actors
with location-abstracted activation: an actor's host is determined by
silo-cluster membership, not by application code. Orleans dissolves the
location-of-the-actor but retains the silo cluster's coordination layer
--- the server-role is internalized in the membership protocol rather
than externalized to the application, but it is not dissolved at the §3
level. Amazon Dynamo (DeCandia et al., 2007) operates as
eventually-consistent without master, with conflict resolution per-key;
it is the canonical AP-without-master system within the CAP positioning
(Brewer, 2000; Gilbert \& Lynch, 2002). Datomic (Cognitect, n.d.), a
commercial immutable database, treats its journal as the source of truth
but routes all writes through a single transactor --- under the §3
reading, the transactor itself is an accidental category that the
present substrate dissolves.

Two theoretical results bracket the structural claim. The CALM theorem
(Hellerstein, 2010; Ameloot et al., 2013) establishes a formal
correspondence between monotonic computation and coordination-freeness
--- a program admits coordination-free distributed execution iff its
computation is monotonic. The substrate of §4 does not appeal to
monotonicity directly, but the dissolution it documents at the role
layer is consistent with CALM's prediction: when state propagation is by
replay of a deterministic program over a journal rather than by
reconciliation of competing writes, the conditions under which
coordination becomes structurally necessary are restricted.
Conflict-free Replicated Data Types (Shapiro et al., 2011), the
mechanism underlying Kleppmann's local-first model, achieve
coordination-free convergence at the data layer; the present substrate
achieves the same property at the program layer, replicating script
invocations rather than data deltas.

The substrate of §4 combines the actor model (Hewitt, 1973), event
sourcing (Fowler, 2005), and CQRS (Young, 2010), synthesized as
practical patterns by Vernon (2015). An additional move --- domain
classes carrying no persistence concerns, discovered by reflection ---
antedates the synthesis by roughly a decade in the framework underlying
the instantiation; that genealogy is described in §4.0.

This paper is the seventh in a series. Papers 1--6 develop the
substrate-level foundations on which the present extension rests,
culminating in Paper 6's \emph{infrastructural symptom}, of which this
paper applies the lens from infrastructural layers to architectural
roles --- a move Paper 6 §9 anticipates explicitly.

The contributions unique to this paper are conceptual and demonstrative.
Conceptually, \emph{accidental category} names the role-level expression
of Paper 6's \emph{infrastructural symptom} --- applying the same two
diagnostic indicators (re-expressed as contingency and dissolution at
the role layer) to the server-role across local-first, peer-to-peer,
agent-centric, and federated systems that prior work treats as piecewise
design targets. Demonstratively, the instantiation of §4 makes the
role-level dissolution operationally visible --- what Paper 6's lens
reads at the role level becomes inspectable in a running three-node
cluster. The §6 observation --- that cloud ecosystems pass from
architectural dependencies to library invocations under the same
conditions --- is a second consequence Paper 6 does not anticipate.

\subsection{3. The construct: accidental category as role-level
infrastructural
symptom}\label{the-construct-accidental-category-as-role-level-infrastructural-symptom}

At the role layer, Paper 6's \emph{infrastructural symptom} typology
reads as \emph{accidental category}. The construct is not new: it is
Paper 6's, applied to a different unit of analysis. Where Paper 6 reads
layers under two diagnostic indicators --- compensation and dissolution
--- this paper re-expresses the same indicators at the role layer as
contingency and dissolution, and applies them to architectural roles.

A server-role reads as an \emph{accidental category} in a class of
systems S when both of two diagnostic indicators are positive.

\textbf{The contingency indicator} (Paper 6's \emph{compensation}
indicator at the role layer). The role's necessity in S is traceable to
a specific representational choice --- what the nodes of S replicate in
order to share state. The choice must be identifiable: not \emph{``the
system uses a server because that is how distributed systems work''},
but \emph{``the system uses a server because the nodes replicate
property R, and the role compensates for the coordination that R
imposes.''} When the contingency indicator is positive, asking
\emph{what would happen if R were replaced?} yields a non-trivial
answer.

\textbf{The dissolution indicator} (Paper 6's \emph{dissolution}
indicator unchanged). When that representational choice is replaced ---
by adopting a substrate in which R is replaced by a different property
--- the role loses justification. The substitution's effect must be
sufficient: nothing else about the problem domain, workload, or
operational context needs to change for the role to become unnecessary.
The constructibility bound of Paper 6 §3 applies here as well: the
substitute substrate must be operationally constructible, not merely an
imagined absence of R.

Both indicators must be positive. Contingency without dissolution
describes a role whose justification survives substrate change --- a
structural requirement of the problem domain. Dissolution without
contingency is not constructible: a role cannot lose justification under
a substrate change unless its justification was tied to the substrate in
the first place. \textbf{The boundary between accidental category and
structural requirement is not a property of the role-as-name but of the
function-for-which-it-is-used.}

The lens is not a universal accusation. Many roles in distributed
systems are structural requirements --- identity providers backed by
cryptographic root authority, gateways mediating external systems,
computationally expensive services such as analytics or inference
engines. No substrate change redistributes the authority of an identity
provider, removes the integration problem of an external network, or
makes intrinsically expensive computation cheap.

The reading operates per role-in-use-case. The same node, identified in
a deployment as ``API server,'' reads as an accidental category when
deployed to accept authoritative writes against a shared database, and
as a structural requirement when deployed to enforce cryptographic
identity claims. The discriminator is the function, not the name.

Seven categories of architectural role are encountered regularly in
production deployments and admit reading under the two indicators. The
third column describes, in substrate-property terms rather than system
terms, what a substrate under which both indicators read positively
offers in place of each role.

\textbf{Table 2.} Diagnostic reading: seven canonical categories of
architectural role under the two indicators of §3. The table identifies,
for each role, the representational choice that necessitates it and the
substrate property that would render it unnecessary.

{\def\LTcaptype{none} % do not increment counter
\begin{longtable}[]{@{}
  >{\raggedright\arraybackslash}p{(\linewidth - 4\tabcolsep) * \real{0.3333}}
  >{\raggedright\arraybackslash}p{(\linewidth - 4\tabcolsep) * \real{0.3333}}
  >{\raggedright\arraybackslash}p{(\linewidth - 4\tabcolsep) * \real{0.3333}}@{}}
\toprule\noalign{}
\begin{minipage}[b]{\linewidth}\raggedright
Architectural role
\end{minipage} & \begin{minipage}[b]{\linewidth}\raggedright
Representational choice that necessitates it
\end{minipage} & \begin{minipage}[b]{\linewidth}\raggedright
What a substrate with both indicators positive offers instead
\end{minipage} \\
\midrule\noalign{}
\endhead
\bottomrule\noalign{}
\endlastfoot
Server as write authority & Nodes replicate data; consistency requires a
single point of acceptance & Each node accepts writes locally; the
journal records the program, which replays deterministically on every
node \\
Master / primary in replication & Nodes replicate state diffs; a primary
must order them & Each node owns its journal; cross-node causality is
recorded by reference, not orchestrated \\
Coordinator in distributed transactions & Multiple writers contend for a
logical entity across processes & Each logical entity has a single point
of authority intrinsic to the substrate; serialization is not
externalized \\
Stateful gateway & Clients cannot directly invoke deferred work without
losing causal record & Deferred work is journal-recorded; the substrate
carries causality, not the gateway \\
Application server / runtime container & The runtime, the framework, and
the application require disparate substrates & A minimal substrate
suffices to run the domain artifact \\
Bootstrap service & New nodes require an issuer of credentials and
addresses & Bootstrap is information transfer that crosses by any means,
including analog artifacts; no service runs \\
Orchestrator as coordination substrate & Stateful coordination across
instances requires external choreography & Coordination is between
actors; the substrate provides location and continuity intrinsically \\
\end{longtable}
}

The seven rows are not exhaustive; service discovery, queue-as-glue, and
configuration management admit similar analysis (§7). The lens is
diagnostic, not prescriptive: it enables auditing existing systems under
the two indicators, independently of any decision to adopt the substrate
of §4.

\subsection{4. Instantiation}\label{instantiation}

\subsubsection{4.0 Genealogy}\label{genealogy}

The substrate underlying the present instantiation is the journaled
actor-native framework whose architectural lineage is documented in
Paper 6 §4.0 and is not retraced here. The implementation source will be
released open-source with the publication of this paper; the references
in this section to specific test methods, file paths, and orchestrator
scripts are pointers into the forthcoming release. Reproduction details
are consolidated in Appendix A.

\subsubsection{4.1 The substrate}\label{the-substrate}

The term \emph{program} throughout this section is used in the
substrate-level sense of Paper 2 §1.2: the pair (domain library, journal
of invocations) that nodes replicate to share state. The substrate
provides three properties that together make the two indicators of §3
read positively for the server-role. First, each node owns a local
journal of invocations --- DSL scripts (in the unit sense of Paper 2)
that, on replay, reconstruct the node's state deterministically. Second,
the domain library --- the classes and verbs that the journal scripts
invoke --- is loaded by reflection at startup; no node hosts the domain
as a service for others. Third, replication is journal-to-journal: nodes
exchange script references and parameters, not state diffs, and reach
convergence by replaying the same sequence of scripts against the same
library. None of these properties requires a node to fulfill the
server-role; each is intrinsic to the substrate, not externalized.

The substrate is described in the abstract throughout this section. Its
concrete instantiation as the framework that this paper's existence
proof builds upon is the one developed across the preceding papers of
the series. The choice of that framework is rhetorical --- it is the one
whose source is available and whose author can speak to its properties
--- not essential to the construct.

\subsubsection{4.2 The port}\label{the-port}

The instantiation ports the \emph{Ordering} bounded context of
Microsoft's \texttt{dotnet/eShop} reference application --- a
widely-used example of microservice-oriented cloud architecture in the
.NET ecosystem --- onto the substrate of §4.1. The diagnostic value of
the port is what it makes visible: the lines of the original codebase
that the substrate cannot host without modification are a small
minority, concentrated in persistence-and-deployment glue. The
aggregate, the value objects, the lifecycle verbs, and the invariants
move into the journaled substrate without modification. What was
previously named ``the Ordering microservice'' turns out to be a clean
domain library wrapped in infrastructure compensating for a different
substrate; the port separates the two.

\subsubsection{4.3 The three-node
deployment}\label{the-three-node-deployment}

Three identical processes --- each running the same substrate with the
ported domain library loaded by reflection --- deploy across three
Docker containers. The coordination mesh is fully connected: three
pairwise channels forming the complete \texttt{N(N–1)/2} graph at
\texttt{N\ =\ 3}. Three is the minimum peer count at which the
dissolution claim is unambiguous: with two peers, any rotation can be
read as one serving the other; with three, no subset of two dominates
the third.

\begin{verbatim}
  ┌─────────────┐    coord pair    ┌─────────────┐
  │ ordering-a  │◄────────────────►│ ordering-b  │
  └──────┬──────┘                  └──────┬──────┘
         │ coord pair                     │
         └──────────┐         ┌───────────┘
                    ▼         ▼
                  ┌─────────────┐
                  │ ordering-c  │
                  └─────────────┘
\end{verbatim}

\textbf{Figure 1.} The three-node coordination mesh. Each node is a
peer; none is privileged.

No node carries the write-acceptance privilege permanently; the role
rotates symmetrically across peers. After the rotation completes, the
three journals are byte-coincident.

The rotation is not consensus-based, in the sense of Raft (Ongaro \&
Ousterhout, 2014) or Multi-Paxos. The substrate does not elect Directors
internally, does not run a voting protocol, does not maintain term or
epoch numbers, and does not provide fencing tokens. Promotion is an
external decision --- the application or operator calls
\texttt{PromoteToDirector()}
(\texttt{Choreography/StageManager/Stage.cs:470}) --- and the new
Director's identity propagates via best-effort announce on the
coordination bus; peers update a local \texttt{directorId} field on
receipt. What the substrate offers is eventual consistency via
deterministic journal replay (Paper 3 develops this as the partition
between work-due-before-responding and deferred work) and cross-actor
causal continuity recorded in the journal rather than orchestrated
externally (Paper 4); it does not offer linearizability, which is a
stronger guarantee that requires consensus and is constructible above
the substrate, not promised by it. If two peers simultaneously believe
themselves to be Director --- a possibility the substrate does not
prevent --- both write to their local journals; on reconnection the
divergence is detectable by journal-hash comparison, and resolution is
per-domain rather than provided by the substrate. The server-role of §3
is the permanent privileged identity that prior topologies require some
node to hold; the substrate dissolves that permanence by making
write-acceptance a transient mode any peer can assume locally. What it
does not remove, and does not claim to remove, is the broader consensus
problem --- a different question with different machinery.

\subsubsection{4.4 The analog bootstrap}\label{the-analog-bootstrap}

The bootstrap of a new peer proceeds in five phases. \textbf{F1} opens
an invitation on the issuer side. \textbf{F2} transmits the invitation
to the joining peer by any means, including analog ones --- printed
paper, a verbal dictation, an email. \textbf{F3} opens an authenticated
connection back to the issuer. \textbf{F4} pauses for a human decision
to approve or reject. \textbf{F5} seals a credential to the joining
peer's public key and closes.

Two of the five phases are sub-software: F2 (the information transfer)
and F4 (the authorization decision). Neither requires a service to be
running anywhere; the carrier of F2 is whatever the operator chooses,
and the decision-maker of F4 is human. The other three phases (F1, F3,
F5) are themselves server-role discharge by the §3 definition: the
issuer listens on an address (F1), accepts a connection (F3), and seals
a credential (F5). The discharge is genuine, not concealed. What §4
documents is that this discharge is ephemeral, time-bounded, and
terminating: the three software-mediated phases are local to the issuer
process, which exits when the bootstrap is complete. The transport stack
that carries F3 (HTTPS/TLS via Kestrel; see Appendix A.4) is a property
of the underlying transport, not of the journaled substrate --- the
substrate's contract is that whatever the transport delivers, the
journal records and replays. \textbf{The issuer process does not act as
a server, coordinator, authority, registry, or rendezvous point after
bootstrap; its only role is to transfer a credential during the brief
lifetime of the handshake, after which it ceases to exist.}

After all three containers complete the handshake --- each with its own
QR code carried by analog means and approved by a human --- the issuer
process exits. The containers continue to operate.

\begin{verbatim}
  ordering-a   Up 11 seconds   0.0.0.0:6443->6443/tcp
  ordering-b   Up 11 seconds   0.0.0.0:6444->6443/tcp
  ordering-c   Up 11 seconds   0.0.0.0:6445->6443/tcp
\end{verbatim}

\textbf{Output of \texttt{docker\ compose\ ps} after the issuer exits.}
No surviving network path connects the containers to the issuer; the
bootstrap papers are in a drawer.

The information that constituted the bootstrap traveled by paper. The
decision that authorized it was made by a human. The cryptographic
protection of the credential happened during the brief lifetime of the
issuer process. The artifact that survives is the three running
containers.

\subsubsection{4.5 The empirical matrix}\label{the-empirical-matrix}

The demo runs across a two-by-two matrix: \emph{in-process} versus
\emph{cross-container}, and \emph{inline} versus \emph{parametric}. Each
cell is a complete execution of the workload across the three nodes; the
in-process row reduces the bootstrap to direct invocation, while the
cross-container row runs the full §4.4 protocol.

\textbf{Table 3.} Empirical matrix of the existence proof. Each cell
reports the final journal entry count, identical on all three peers and
byte-coincident on disk.

{\def\LTcaptype{none} % do not increment counter
\begin{longtable}[]{@{}
  >{\raggedright\arraybackslash}p{(\linewidth - 4\tabcolsep) * \real{0.3333}}
  >{\raggedright\arraybackslash}p{(\linewidth - 4\tabcolsep) * \real{0.3333}}
  >{\raggedright\arraybackslash}p{(\linewidth - 4\tabcolsep) * \real{0.3333}}@{}}
\toprule\noalign{}
\begin{minipage}[b]{\linewidth}\raggedright
Workload regime
\end{minipage} & \begin{minipage}[b]{\linewidth}\raggedright
In-process (3 Stages, single OS process)
\end{minipage} & \begin{minipage}[b]{\linewidth}\raggedright
Cross-container (3 Dockers, real TLS + analog bootstrap)
\end{minipage} \\
\midrule\noalign{}
\endhead
\bottomrule\noalign{}
\endlastfoot
Inline (literal scripts) & 22 entries & 22 entries \\
Parametric (Define + bind) & 7 entries & 7 entries \\
\end{longtable}
}

The four cells produce identical final entry counts per workload regime,
and the three peers in each cell are byte-coincident on disk. The
cross-container row confirms that the operational properties of the
substrate survive the addition of the analog bootstrap and the real
cryptographic stack; the in-process row is the structural rehearsal that
isolates the substrate's intrinsic properties from the network. The
compaction from inline to parametric is real and measured; its precise
ratio and the conditions under which it varies are treated in Paper 2.

Byte-coincidence is the structural witness of the §3 reading: it shows
that no peer is privileged --- the four functions of the server-role
definition (write authority, coordination ownership, mediation,
addressing) are met symmetrically across peers rather than by a single
node. Operational metrics --- write throughput, append latency, cross-DC
convergence --- live elsewhere in the series: Paper 5 lab L1 reports
451k events/sec on the substrate, lab L6 reports 347µs p50 append on the
FS backend, and lab L4 demonstrates bit-exact cross-DC convergence under
sustained workload. The present experiment is not a benchmark; it is the
existence proof for the role-level reading of Paper 6's lens, alongside
the operational floor that Paper 5 establishes.

In the parametric regime, the journal records seven entries rather than
five: each rotation Director emits its own Define + Invocation pair,
because the Define cache is local to the Stage that emits it, not
replicated. The journal is the catalog of each peer's declarative
vocabulary; the cache's compactness lives on the Stage axis, not on the
journal axis.

\subsubsection{4.6 Partition tolerance and
re-join}\label{partition-tolerance-and-re-join}

The same artifact additionally exhibits partition tolerance and re-join
from disk. A peer can exit the cluster while the others continue
operating; a fresh process can mount the absent peer's data directory
and rehydrate its journal-relative state from disk alone, without
contacting the issuer or the surviving peers; and the rehydrated peer
integrates into the still-running cluster via peer-to-peer catch-up. The
five-phase decomposition of this scenario and its implications for the
construct's dissolution claim are described in §5. The relevant
observation for §4 is that the same artifact that produces the four
cells of Table 3 also produces, without modification, the operational
signature of partition tolerance and re-join from disk.

\subsection{5. The first consequence: dissolution of the
server-role}\label{the-first-consequence-dissolution-of-the-server-role}

§3 defined the server-role as a composite of four architectural
functions. §4 instantiated a substrate in which each of these functions
is distributed intrinsically rather than externalized to a privileged
node. This section reports the dissolution as it occurs across the
operational lifecycle of the cluster.

Acceptance of authoritative writes and ownership of coordination
redistribute together under Director rotation. The Director role
traverses the three peers in turn; no node carries the write-acceptance
privilege permanently; the privilege is a transient mode that any peer
can assume. The same property holds across the workload's non-happy
paths --- cancellation, error recovery --- not only across its main
sequence.

Mediation between clients and provision of the cluster's address
distribute into the peer-to-peer mesh. No node carries an authoritative
address for the cluster as a whole; a request directed at any peer is
acceptable and that peer can act as Director for it. The ``address
against which other parties direct their traffic'' --- the fourth limb
of the §3 definition --- is not a single address but the union of three
pinned addresses, none of which is privileged.

The strongest empirical signature of the role's dissolution is the
partition-tolerance scenario introduced in §4.6. One peer leaves the
cluster; the remaining two continue to operate; the absent peer rejoins
by reading its on-disk journal. Five phases describe the sequence:

\begin{itemize}
\tightlist
\item
  \textbf{R1.} The three peers converge to a shared journal state.
\item
  \textbf{R2.} One peer exits; the others continue. The cluster retains
  quorum, retains a Director, and accepts new work.
\item
  \textbf{R3.} The surviving Director issues additional work; the two
  surviving journals advance; the absent peer's on-disk journal is
  unchanged.
\item
  \textbf{R4.} A fresh process --- distinct in-memory identity, same
  data directory --- reads the disk and reconstructs the
  journal-relative state of the absent peer, without contacting the
  issuer and without contacting the surviving peers.
\item
  \textbf{R5.} The rehydrated peer integrates into the still-running
  cluster via peer-to-peer catch-up; the gap is replayed from a
  surviving peer; subsequent work brings all three peers back to a
  shared state.
\end{itemize}

The properties exercised by this scenario are not separable. The
cluster's continued operation in the absent peer's absence is only
possible if no node is the cluster's source of authority. The
rehydration without consulting the issuer is only possible if the disk
is a self-contained witness of the peer's prior participation. The
catch-up between peers is only possible if any surviving peer is capable
of replaying the gap; no specialized recovery node is invoked.
\emph{Operationally, the scenario shows that a peer can leave the
cluster, the cluster can continue accepting work in its absence, and the
peer can re-join by reading the journal that survived its disappearance
on its own disk --- no central authority is consulted at any point in
this loop. The role of the issuer was discharged once, at the original
onboarding, and is not invoked again.}

The issuer is invoked exactly once per peer. The credential sealed
during the bootstrap of §4.4 is preserved across restarts, and the
peer's continued participation is verified by the journal's
cryptographic integrity, not by reference to a running service. If a
peer loses its disk, it re-onboards as a new peer; if it preserves its
disk, it returns as itself, regardless of how long it was absent.

The four functions of the §3 definition, taken together, are what the
literature names ``server.'' The §4 substrate redistributes them in
turn: write acceptance and coordination ownership distribute across the
three peers via rotation; mediation and addressing distribute across the
mesh; partition tolerance and re-join distribute across the surviving
peers via catch-up; the bootstrap issuer participates once and is absent
thereafter. \textbf{The server is not a permanent identity. The server
is a transferable coordinator role that any peer can assume in turn ---
and any peer can leave the cluster and re-join without reference to a
coordinator outside it.}

\subsection{6. The second consequence: cloud as library, not
infrastructure}\label{the-second-consequence-cloud-as-library-not-infrastructure}

A second operational consequence becomes visible in §4 once the cluster
continues to operate after bootstrap without any of the services that
originally constituted the Ordering microservice. Cloud and microservice
ecosystems, when accessed from a journaled program, are observable as
libraries invoked by the program rather than as services required for
the cluster to exist.

The structural difference is not in what happens when an external
service is unreachable --- resilience patterns developed for
microservices already achieve graceful degradation under that condition
--- but in the substrate's relation to external invocation by default.
Resilience patterns (Hystrix at Netflix; Polly in .NET; sidecar-based
circuit breakers in Istio; the catalog in Nygard, 2018) achieve
library-like consumption at the application level: each external call is
wrapped in an opt-in policy, each fallback handler is explicitly
written, and the discipline of doing this consistently is a property of
the development team, not of the runtime. Under the journaled substrate,
library-like consumption is the default: every external invocation is
recorded as a journal entry whether or not the developer thought to wrap
it, the outcome (including failure) is part of the program's durable
state, and the journal is replayable as part of the program's history.
The substrate makes the property structural rather than disciplinary;
the journal-recorded failure is not merely observability metadata but
program state available to subsequent Reactions (Paper 3) and Tell
continuations (Paper 4).

After the port described in §4.2, none of the services that previously
had to be running for the Ordering microservice to exist are reachable
from the cluster. The domain code runs entirely inside each node; what
previously required an API host, a database, and an event bus is now a
library invoked by journaled programs. The cluster operates against its
own journals; the cloud topology of the original microservice has no
point of contact with the cluster after the bootstrap of §4.4 completes.

The same relation generalizes to external services that the cluster
invokes during normal operation. If the identity provider is unreachable
at the moment of invocation, the journal records the failed invocation
and the cluster continues. If a payment gateway is unreachable, the
journal records the attempt and its outcome. The external service's
absence is recorded as data, not experienced as outage. The cluster does
not inhabit any of these services; it invokes them at chosen moments
through journaled Reactions (Paper 3) and proceeds with their result,
including the result of their absence.

This operational relation inverts the formulation that has organized the
local-first software literature:

\textbf{Table 4.} The inversion of Kleppmann's formulation, captured as
a six-word pair.

{\def\LTcaptype{none} % do not increment counter
\begin{longtable}[]{@{}
  >{\raggedright\arraybackslash}p{(\linewidth - 2\tabcolsep) * \real{0.5000}}
  >{\raggedright\arraybackslash}p{(\linewidth - 2\tabcolsep) * \real{0.5000}}@{}}
\toprule\noalign{}
\begin{minipage}[b]{\linewidth}\raggedright
Kleppmann (Local-first)
\end{minipage} & \begin{minipage}[b]{\linewidth}\raggedright
The present construct
\end{minipage} \\
\midrule\noalign{}
\endhead
\bottomrule\noalign{}
\endlastfoot
\emph{Own your data in spite of the cloud.} & \emph{Record the cloud;
don't inhabit it.} \\
\end{longtable}
}

Local-first software treats the cloud as something the software must
survive without. The cluster of §4 neither survives without the cloud
nor depends on it; it invokes cloud services when useful and proceeds
without them when not. The two relations differ in what the absence of a
cloud service means for the cluster's continued existence.

The experiment shows that cloud services, when used from a journaled
program, are consumed as invocable capabilities rather than as the
substrate in which the program must reside.

\subsection{7. Limitations and
counter-arguments}\label{limitations-and-counter-arguments}

The instantiation of §4 carries specific limitations that follow from
the choices made for the demo. They are reported here honestly; none of
them affects the lens's reading, but each affects what the present
experiment alone establishes.

\textbf{Identity in the port is local.} The original
\texttt{dotnet/eShop} deployment uses an external identity provider; the
port substitutes a local string-typed identifier. The lens does not read
identity providers as dissolving under the indicators of §3 --- identity
providers backed by cryptographic root authority are structural
requirements, as Table 2 notes --- and the demo does not exhibit one. A
second experiment that integrates an external identity provider as a
Reaction-invoked library would extend the empirical surface without
changing the lens's reading.

\textbf{Cross-context dependencies are not portrayed.} The Ordering
bounded context of \texttt{dotnet/eShop} carries references resolved by
the Catalog and Basket bounded contexts; the port treats those
references as opaque identifiers. The lens reads the same port mechanism
as applying to Catalog and Basket; the present experiment exhibits
Ordering only.

The two limitations above warrant clarification, because they are easy
to read as selection bias when they are in fact bounded-context
isolation in the Domain-Driven Design sense (Evans, 2003; Vernon, 2013).
The Ordering bounded context, in any architecture --- microservices,
monolith, or journaled-program substrate --- references Catalog and
Basket entities by identifier rather than by full object, because
cross-context object resolution is itself an integration boundary, not a
property of Ordering's domain. Identity, similarly, lives in its own
bounded context (the Identity microservice in the original eShop);
Ordering does not authenticate users, it accepts a \texttt{UserId}.
Removing those from the port is not eliminating components that ``need a
server''; it is respecting the same boundaries DDD prescribes. Table 2
itself lists identity providers backed by cryptographic root authority
as structural requirements --- the lens never reads identity as
dissolving, so its exclusion is not adverse evidence. What the port
demonstrates is server-role dissolution for the Ordering bounded
context; the same lens applied to Catalog or Basket would be a separate
exercise, and the same lens applied to Identity would read it as
structural.

\textbf{The empirical matrix is at N = 3.} The mesh in §4.3 connects
three peers --- the minimum non-trivial fully-connected mesh. The lens's
reading is not bounded at N = 3; the indicators of §3 are
topology-agnostic and node-count-agnostic. The fully-connected mesh of
§4.3 is a demo choice --- the simplest topology in which write-ownership
rotation is unambiguously visible at the minimum non-trivial peer count.
At larger N the operator chooses an overlay appropriate to scale
(gossip, hierarchical overlays, sharding) --- the same techniques by
which Bluesky's relay network operates at federation scale --- and the
indicators read positively under any such overlay that preserves
journal-determinism between peers.

\textbf{The cross-container partition-tolerance scenario is captured in
a parallel artefact, not in the matrix of §4.} The natural
cross-container equivalent of the in-process scenario in §4.6 ---
stopping a container, allowing the running cluster to continue accepting
work in its absence, restarting the container, and watching it rehydrate
from its mounted volume and rejoin via catch-up against a surviving peer
--- is implemented in the same orchestrator that produced the matrix
(the \texttt{-\/-rehydrate-demo} flag of the demo script) and depends on
no substrate feature that is not already exercised by the in-process
test in §4.6. The decision to keep §4 to the four cells of the matrix,
rather than expanding to a fifth cell, was rhetorical: the in-process
exercise establishes the property structurally, and §5 reasons from
there. The cross-container capture documents the same property under a
different process boundary; it does not add a new claim.

\textbf{The compaction ratio between inline and parametric workloads
depends on Director rotation.} Each new Director journals a fresh Define
+ Invocation pair, so the Define cache compacts on the Stage axis rather
than on the journal axis. The compaction is real and measured at three
rounds; its precise value at larger workloads is treated in Paper 2.

The lens's reading extends to other architectural roles that admit the
same analysis. Service discovery, queue-as-glue between bounded
contexts, configuration-management roles, and distributed-tracing
collectors are auditable under the two indicators of §3; the present
experiment does not exercise them, and the substrate's relation to each
is left as a forward question. The lens admits this open extension by
design --- its diagnostic value is in reading which categories in any
given system are contingent, not in enumerating them exhaustively.

\subsection{8. Conclusion}\label{conclusion}

This paper has extended Paper 6's \emph{infrastructural symptom} lens
from infrastructural layers to architectural roles, applying it to a
property that was previously implicit in distributed-systems practice:
that the server-role is not necessarily structural to the problem domain
but can arise from representational choices in the persistence model.
The term \emph{accidental category} names this role-level expression.

The two diagnostic indicators of Paper 6 --- re-expressed as contingency
and dissolution at the role layer --- allow practitioners to read
accidental categories apart from structural requirements, and §3 applies
them across canonical architectural roles. The instantiation of §4
demonstrates that, under a journaled-program substrate, the Ordering
bounded context of \texttt{dotnet/eShop} continues to operate across
three mutually bootstrapping nodes without any participant fulfilling
the server-role. The demonstration does not argue uniqueness of the
substrate; it establishes that the role-level reading of Paper 6's lens
is operationally inspectable.

Under these conditions, the four functions traditionally attributed to
the server-role --- authoritative writes, coordination ownership, client
mediation, and address provision --- are observed to become properties
of the substrate itself. What appears as ``server behavior'' is
redistributed across peer rotation, mesh communication, disk-backed
partition tolerance, and one-time bootstrap discharge.

The same observation extends to external infrastructure. Cloud services,
microservices, and infrastructural ecosystems, when accessed from a
journaled program, are no longer required habitats for the system to
exist, but libraries invoked by the program when needed. The distinction
is not eliminative but relational: from inhabitation to invocation. This
second consequence is novel to this paper --- Paper 6 does not
anticipate it.

The contributions of this paper are two. Conceptually, the role-level
extension of Paper 6's lens can be applied independently of the
substrate used in §4 --- any production system can be examined
role-by-role, asking whether each role is tied to the problem being
solved or to the persistence model chosen. Operationally, the working
instantiation makes that examination concretely inspectable in one
well-known reference case.

The reading is per role-in-use-case, not per system. Identity
authorities, payment networks, sensor integrations, and computationally
intensive services remain structural under any substrate. The lens does
not promise universal dissolution; it offers examinability.

Paper 6 applied the lens at the layer level, identifying many
infrastructural layers --- caches, queues, locks, application servers,
orchestration clusters --- as compensations for the underlying
persistence model. The present paper applies the same lens to the
architectural role that made such layers necessary in the first place:
the server. The two papers apply one lens at two levels of abstraction
--- first to layers, then to roles. From this point on, treating the
server-role as a structural requirement is no longer an assumption but a
choice that can be made explicit and evaluated.

\subsection{Appendix A ---
Reproducibility}\label{appendix-a-reproducibility}

The artifact described in §4 is publicly available. The references in
this appendix to file paths, test method names, and orchestrator scripts
resolve against the released code; the artifact comprises the substrate
framework, the runtime, the ported \texttt{Ordering} domain library and
its test suite, the demo orchestrator, the supporting reports, and the
asciinema recordings, hosted across two repositories described in §A.1.

\subsubsection{A.1 Source and suite}\label{a.1-source-and-suite}

The artifact is hosted in two public repositories under
https://github.com/alvaroNCubo. The substrate framework
(\texttt{Choreography/}), the runtime (\texttt{Puppeteer/}), the ported
\texttt{Ordering} domain library and its test suite
(\texttt{tests-local/UnitTestPaper7EShop/}), the demo orchestrator
(\texttt{docker/}), and the supporting reports referenced in §A.7
(\texttt{notes/paper7\_phase1\_results.md},
\texttt{notes/paper7\_phase2\_results.md}) are in
https://github.com/alvaroNCubo/puppeteer. The asciinema recordings and
GIFs of §A.8 are in the \texttt{paper7-assets/} directory of
https://github.com/alvaroNCubo/puppeteer-papers. Inline code citations
of the form \texttt{file.cs:NN} resolve against the commit pinned in
Code Provenance.

\subsubsection{A.2 Demo orchestrator}\label{a.2-demo-orchestrator}

The orchestrator script \texttt{docker/run-demo.sh} provisions three
Docker containers, runs the analog bootstrap handshake against each, and
exercises one cell of the matrix per invocation. Three flags select the
scenario:

\begin{itemize}
\tightlist
\item
  \texttt{bash\ docker/run-demo.sh} --- cross-container inline (22
  entries final).
\item
  \texttt{bash\ docker/run-demo.sh\ -\/-parametric} --- cross-container
  parametric (7 entries final).
\item
  \texttt{bash\ docker/run-demo.sh\ -\/-rehydrate-demo} --- captures the
  cross-container leave-and-rejoin scenario referenced in §7: stops one
  container, allows the cluster to advance, restarts the container, and
  observes its rehydration and catch-up.
\end{itemize}

Each invocation runs in approximately 60--90 seconds on Docker Desktop.

\subsubsection{A.3 Tests cited}\label{a.3-tests-cited}

The empirical claims in the body are exercised by specific tests in the
suite:

{\def\LTcaptype{none} % do not increment counter
\begin{longtable}[]{@{}
  >{\raggedright\arraybackslash}p{(\linewidth - 6\tabcolsep) * \real{0.2500}}
  >{\raggedright\arraybackslash}p{(\linewidth - 6\tabcolsep) * \real{0.2500}}
  >{\raggedright\arraybackslash}p{(\linewidth - 6\tabcolsep) * \real{0.2500}}
  >{\raggedright\arraybackslash}p{(\linewidth - 6\tabcolsep) * \real{0.2500}}@{}}
\toprule\noalign{}
\begin{minipage}[b]{\linewidth}\raggedright
Body section
\end{minipage} & \begin{minipage}[b]{\linewidth}\raggedright
Property
\end{minipage} & \begin{minipage}[b]{\linewidth}\raggedright
Test method
\end{minipage} & \begin{minipage}[b]{\linewidth}\raggedright
Project
\end{minipage} \\
\midrule\noalign{}
\endhead
\bottomrule\noalign{}
\endlastfoot
§4.4 & F1--F5 bootstrap handshake under real cryptography &
\texttt{EndToEnd\_RealCryptoOverRealTls\_RoundsTripIdentity} &
\texttt{UnitTestChoreography/UsherOnboardingTests.cs} \\
§5 (single peer operates alone) & Force-promoted Stage operates without
peers & \texttt{SingleStage\_OperatesAlone\_WhenPromotedForce} &
\texttt{UnitTestPaper7EShop/ThreeNodeOrderingTests.cs} \\
§4.5 (in-process, inline) & Three-stage Director rotation, inline regime
& \texttt{ThreeStages\_DirectorRotation\_AllConverge\_HappyPath} &
\texttt{UnitTestPaper7EShop/ThreeNodeOrderingTests.cs} \\
§4.5 (in-process, parametric) & Three-stage Director rotation,
parametric regime &
\texttt{ThreeStages\_DirectorRotation\_AllConverge\_HappyPath\_Parametric}
& \texttt{UnitTestPaper7EShop/ThreeNodeOrderingTests.cs} \\
§5 (rotation under cancellation) & Cancellation branch replicates across
three stages &
\texttt{CancellationBranch\_Replicates\_AcrossThreeStages} &
\texttt{UnitTestPaper7EShop/ThreeNodeOrderingTests.cs} \\
§4.5 (instrumentation) & DLL load + per-phase convergence + journal
bytes snapshot & \texttt{Metrics\_F1\_Snapshot} &
\texttt{UnitTestPaper7EShop/ThreeNodeOrderingTests.cs} \\
§4.6 / §5 (partition tolerance) & Leave-and-rejoin from disk,
peer-to-peer catch-up &
\texttt{OneStage\_Offline\_OthersAdvance\_Rehydrate\_CatchUp\_Parametric}
& \texttt{UnitTestPaper7EShop/ThreeNodeOrderingTests.cs} \\
§5 (rehydration without onboarding) & Rehydrated Stage restores identity
from disk alone &
\texttt{RehydratedStage\_RestoresIdentityAndJournal\_WithoutOnboarding}
& \texttt{UnitTestPaper7EShop/ThreeNodeOrderingTests.cs} \\
\end{longtable}
}

\subsubsection{A.4 Cryptographic stack}\label{a.4-cryptographic-stack}

The cryptographic primitives underlying the F1--F5 handshake of §4.4 are
documented in \texttt{Choreography/Usher/Crypto/*.cs}. The
implementations use BouncyCastle 2.6.2:

\begin{itemize}
\tightlist
\item
  Ed25519 keypair generation, signing, verification:
  \texttt{Ed25519StageKeyGenerator.cs}, \texttt{Ed25519Signer.cs},
  \texttt{Ed25519SignatureVerifier.cs}.
\item
  Ed25519-to-X25519 derivation (Edwards-to-Montgomery, \emph{u = (1 + y)
  / (1 − y) mod p}): \texttt{Ed25519ToX25519.cs}.
\item
  Sealed-box AEAD using X25519 ECDH and ChaCha20-Poly1305, with the
  participating public keys bound as additional authenticated data:
  \texttt{SealedBoxPayloadSealer.cs}.
\item
  Kestrel TLS listener and self-signed certificate generation with SAN:
  \texttt{HttpsTransportListener.cs}, \texttt{SelfSignedCert.cs}.
\item
  TOFU certificate fingerprint pin via
  \texttt{SocketsHttpHandler.SslOptions.RemoteCertificateValidationCallback}:
  \texttt{HttpsClientFactory.cs}.
\end{itemize}

\subsubsection{A.5 Share-link URI}\label{a.5-share-link-uri}

The share-link emitted by the issuer takes the form
\texttt{puppeteer-usher://v1/\{base64url(json)\}}. The JSON payload
carries the transport address of the issuer endpoint and a SHA-256
fingerprint (\texttt{fp} field) of its TLS certificate. The
encoder/decoder is defined by \texttt{IUsherShareLinkEncoder} in
\texttt{Choreography/Usher/}. A representative share-link under load is
449 characters in length, within the alphanumeric capacity of a QR code
at version 14 with error correction level M (approximately 535
characters).

\subsubsection{A.6 Wire format}\label{a.6-wire-format}

The HTTPS transport between containers carries a \texttt{ChannelPurpose}
byte after the sender identifier in each frame. This wire format (v2)
prevents Define and Invocation channels from collapsing into the same
channel-dictionary key when Director rotation runs the parametric
workload across multiple containers; the in-process regime is
unaffected. The wire format is implemented in the transport handlers
under \texttt{Choreography/Transport/Https/}.

\subsubsection{A.7 Supporting reports}\label{a.7-supporting-reports}

Two markdown reports in the \texttt{notes/} directory of the same branch
document the experiment in detail:

\begin{itemize}
\tightlist
\item
  \texttt{notes/paper7\_phase1\_results.md} --- the in-process row of
  the matrix (eleven sections plus three addenda).
\item
  \texttt{notes/paper7\_phase2\_results.md} --- the cross-container row,
  the hardening cycle that converged on the demo, the three-Docker mesh,
  the rotation protocol, and the parametric cell.
\end{itemize}

\subsubsection{A.8 Recordings}\label{a.8-recordings}

Eleven cast recordings (asciinema format) and eleven GIF renders capture
the four cells of the matrix of §4.5 and the rehydration property of
§4.6, one recording per Docker container per cell. The recordings are
published alongside this paper as supplementary material, hosted in the
\texttt{paper7-assets/} directory of the \texttt{puppeteer-papers}
repository.

\subsection{Code provenance}\label{code-provenance}

Source-code references in this paper resolve against the public
Puppeteer repository at commit
\href{https://github.com/alvaroNCubo/puppeteer/tree/91118fc1a3d865e0bb44e0490712c46ceacf5c63}{\texttt{91118fc}}
(2026-05-22). The substrate code cited inline (file:line references in
§4-§5) is identical to the earlier snapshot at commit \texttt{2f31f96}
(2026-05-18); commit \texttt{91118fc} only adds the reproducibility
artifacts of Appendix A (\texttt{docker/}, \texttt{notes/}). The earlier
substrate snapshot is archived in Software Heritage under the following
persistent identifier:

\begin{verbatim}
swh:1:dir:10e7e6bad7eb77b6c2e406762026177f95c3ae92;
  origin=https://github.com/alvaroNCubo/puppeteer;
  anchor=swh:1:rev:2f31f9674a5de816bdf1bf9d8360ff218a02e4da
\end{verbatim}

Inline references of the form \texttt{file.cs:NN} (e.g.,
\texttt{ActorHandler.cs:38}) resolve against this snapshot. A reader can
construct a per-file SWHID by adding the qualifiers
\texttt{;path=\textless{}path\textgreater{};lines=\textless{}NN\textgreater{}}
to the directory SWHID above. Future commits to the repository may
renumber lines; the SWHID preserves the cited state independently of any
future change to the repository or its hosting.

\subsection{Acknowledgments}\label{acknowledgments}

The author used large language models (including Claude and ChatGPT) as
editorial assistants for language refinement, structural feedback, and
literature navigation. All original ideas, terminology, theoretical
constructs, and technical content presented in this work are solely the
author's.

\subsection{References}\label{references}

Ameloot, T. J., Neven, F., \& Van den Bussche, J. (2013). Relational
transducers for declarative networking. \emph{Journal of the ACM},
60(2), 1--38.

Armstrong, J. (2003). \emph{Making reliable distributed systems in the
presence of software errors} {[}Doctoral thesis{]}. KTH Royal Institute
of Technology.

Bernstein, P. A., Bykov, S., Geller, A., Kliot, G., \& Thelin, J.
(2014). \emph{Orleans: Distributed virtual actors for programmability
and scalability} (Tech. Rep.~MSR-TR-2014-41). Microsoft Research.

Brewer, E. A. (2000). Towards robust distributed systems {[}Keynote
address{]}. \emph{Proceedings of the Nineteenth Annual ACM Symposium on
Principles of Distributed Computing (PODC '00)}.

Cognitect. (n.d.). \emph{Datomic: A transactional database with a
flexible data model, queryable as a log of changes} {[}Software
project{]}. Retrieved May 22, 2026, from https://www.datomic.com/

DeCandia, G., Hastorun, D., Jampani, M., Kakulapati, G., Lakshman, A.,
Pilchin, A., Sivasubramanian, S., Vosshall, P., \& Vogels, W. (2007).
Dynamo: Amazon's highly available key-value store. \emph{Proceedings of
the 21st ACM SIGOPS Symposium on Operating Systems Principles (SOSP
'07)}, 205--220.

Earthstar Project. (n.d.). \emph{Earthstar: simple personal data sync}
{[}Software project{]}. Retrieved May 22, 2026, from
https://earthstar-project.org/

Evans, E. (2003). \emph{Domain-driven design: Tackling complexity in the
heart of software}. Addison-Wesley.

Fowler, M. (2005). Event sourcing. Retrieved from
https://martinfowler.com/eaaDev/EventSourcing.html

Gilbert, S., \& Lynch, N. (2002). Brewer's conjecture and the
feasibility of consistent, available, partition-tolerant web services.
\emph{ACM SIGACT News}, 33(2), 51--59.

Gregor, S. (2006). The nature of theory in information systems.
\emph{MIS Quarterly}, 30(3), 611--642.

Harris-Braun, E., Luck, N., \& Brock, A. (2018). \emph{Holochain:
Scalable agent-centric distributed computing} (Alpha 1 white paper).
Holo.
https://www.holochain.org/documents/holochain-white-paper-alpha.pdf

Hellerstein, J. M. (2010). The declarative imperative: Experiences and
conjectures in distributed logic. \emph{ACM SIGMOD Record}, 39(1),
5--19.

Hewitt, C., Bishop, P., \& Steiger, R. (1973). A universal modular ACTOR
formalism for artificial intelligence. \emph{Proceedings of the 3rd
International Joint Conference on Artificial Intelligence (IJCAI)},
235--245.

Holepunch. (n.d.). \emph{Hypercore: append-only logs for the
peer-to-peer web} {[}Software project{]}. Retrieved May 22, 2026, from
https://docs.pears.com/building-blocks/hypercore

Kleppmann, M., Frazee, P., Gold, J., Graber, J., Holmgren, D., Ivy, D.,
Johnson, J., Newbold, B., \& Volpert, J. (2024). Bluesky and the AT
Protocol: Usable decentralized social media. In \emph{Proceedings of the
ACM CoNEXT-2024 Workshop on the Decentralization of the Internet}. ACM.
https://doi.org/10.1145/3694809.3700740

Kleppmann, M., Wiggins, A., van Hardenberg, P., \& McGranaghan, M.
(2019). Local-first software: You own your data, in spite of the cloud.
In \emph{Onward! 2019: Proceedings of the 2019 ACM SIGPLAN International
Symposium on New Ideas, New Paradigms, and Reflections on Programming
and Software} (pp.~154--178). ACM.
https://doi.org/10.1145/3359591.3359737

Lemmer-Webber, C., Tallon, J., Shepherd, E., Guy, A., \& Prodromou, E.
(2018). \emph{ActivityPub} (W3C Recommendation, 23 January 2018). World
Wide Web Consortium. https://www.w3.org/TR/activitypub/

Matrix.org Foundation. (n.d.). \emph{Matrix specification} {[}Online
document{]}. Retrieved May 22, 2026, from https://spec.matrix.org/

Microsoft. (2026). \emph{dotnet/eShop: A reference .NET application
implementing an eCommerce site} {[}Software repository, commit 9b4f943,
2026-04-21{]}. Retrieved May 22, 2026, from
https://github.com/dotnet/eShop

Nygard, M. T. (2018). \emph{Release it! Design and deploy
production-ready software} (2nd ed.). Pragmatic Bookshelf.

Ogden, M., McKelvey, K., \& Madsen, M. B. (2017). \emph{Dat: Distributed
dataset synchronization and versioning} {[}White paper{]}. Code for
Science. https://github.com/datprotocol/whitepaper

Ongaro, D., \& Ousterhout, J. (2014). In search of an understandable
consensus algorithm. \emph{Proceedings of the 2014 USENIX Annual
Technical Conference (USENIX ATC '14)}, 305--319.

Rivera, A. (2026a). Anti-porous architecture: a unified design principle
for CQRS + Actor + Event-Sourcing systems. \emph{Puppeteer Papers
Series}, Paper 1.
https://github.com/alvaroNCubo/puppeteer-papers/blob/main/01-anti-porosity.md

Rivera, A. (2026b). Program--value separability: the structural
precondition for compilation, caching, and dense journaling in a DSL
runtime. \emph{Puppeteer Papers Series}, Paper 2.
https://github.com/alvaroNCubo/puppeteer-papers/blob/main/02-program-value-separability.md

Rivera, A. (2026c). Reactions and the partition: opt-in eventual
consistency in actor-native systems. \emph{Puppeteer Papers Series},
Paper 3.
https://github.com/alvaroNCubo/puppeteer-papers/blob/main/03-reactions-and-partition.md

Rivera, A. (2026d). Preserving semantic continuity across actors: a
tell-based approach without orchestration. \emph{Puppeteer Papers
Series}, Paper 4.
https://github.com/alvaroNCubo/puppeteer-papers/blob/main/04-cross-actor-continuity.md

Rivera, A. (2026e). The journal as substrate: unifying deployment,
replication, backup, and offline operation in distributed systems.
\emph{Puppeteer Papers Series}, Paper 5.
https://github.com/alvaroNCubo/puppeteer-papers/blob/main/05-substrate-operations.md

Rivera, A. (2026f). Most infrastructure layers are symptoms of the
persistence model: a construct for auditing production stacks.
\emph{Puppeteer Papers Series}, Paper 6.
https://github.com/alvaroNCubo/puppeteer-papers/blob/main/06-infrastructural-symptom.md

Sambra, A. V., Mansour, E., Hawke, S., Zereba, M., Greco, N., Ghanem,
A., Zagidulin, D., Aboulnaga, A., \& Berners-Lee, T. (2016).
\emph{Solid: A platform for decentralized social applications based on
linked data} (Technical Report). MIT CSAIL \& Qatar Computing Research
Institute. http://emansour.com/research/lusail/solid\_protocols.pdf

Shapiro, M., Preguiça, N., Baquero, C., \& Zawirski, M. (2011).
\emph{Conflict-free replicated data types} (Research Report RR-7687).
INRIA.

Tarr, D., Lavoie, E., Meyer, A., \& Tschudin, C. (2019). Secure
Scuttlebutt: An identity-centric protocol for subjective and
decentralized applications. In \emph{Proceedings of the 6th ACM
Conference on Information-Centric Networking (ICN '19)} (pp.~1--11).
ACM. https://doi.org/10.1145/3357150.3357396

Vernon, V. (2013). \emph{Implementing domain-driven design}.
Addison-Wesley.

Vernon, V. (2015). \emph{Reactive messaging patterns with the actor
model: Applications and integration in Scala and Akka}. Addison-Wesley.

Young, G. (2010). \emph{CQRS documents}. Retrieved from
https://cqrs.files.wordpress.com/2010/11/cqrs\_documents.pdf

\end{document}
